\documentclass[11pt]{amsart}
\usepackage{amsmath,amssymb}
\usepackage[colorlinks=true,linkcolor=blue]{hyperref}
\usepackage[nameinlink,noabbrev]{cleveref}

\theoremstyle{plain}
\newtheorem{theorem}{Theorem}
\newtheorem{lemma}[theorem]{Lemma}

\title{Convergence of a Discrete Gradient Flow}
\author{S. Park}
\date{}

\begin{document}
\begin{abstract}
We prove that the implicit Euler discretization of a gradient flow for a
strongly convex potential converges linearly to the unique minimizer.
The main estimate is \cref{thm:linear} in \cref{sec:analysis}, which rests on
the descent lemma \cref{lem:descent} and is illustrated by the flow diagram
\cref{fig:flow} of \cref{sec:method}.
\end{abstract}
\maketitle

\section{Introduction}
\label{sec:intro}

Let $f \colon \mathbf{R}^n \to \mathbf{R}$ be twice differentiable,
$\mu$-strongly convex, and $L$-smooth. The continuous gradient flow
$\dot x = -\nabla f(x)$ converges exponentially to the minimizer
$x^\star$, and our goal is a matching discrete theory. The scheme is defined
in \cref{sec:method}; its convergence is proved in \cref{sec:analysis}.
\Cref{thm:linear} gives the headline linear rate, while
\cref{eq:implicit} displays the update and \cref{fig:flow} sketches one
iteration. Forward pointers of this kind stress the cross-pass machinery:
none of \cref{sec:method}, \cref{sec:analysis}, \cref{eq:implicit},
\cref{lem:descent}, \cref{thm:linear}, or \cref{fig:flow} is defined yet at
the point it is first cited.

\section{Method}
\label{sec:method}

Given a step size $\tau > 0$ and an iterate $x_k$, the implicit Euler step
solves
\begin{equation}
  \label{eq:implicit}
  x_{k+1} = x_k - \tau \nabla f(x_{k+1}),
\end{equation}
equivalently $x_{k+1} = \operatorname{prox}_{\tau f}(x_k)$. One step of
\cref{eq:implicit} is sketched in \cref{fig:flow}; the analysis that turns
this picture into a rate is \cref{thm:linear} in \cref{sec:analysis}.

\begin{figure}[ht]
  \centering
  \rule{0.7\textwidth}{2.5cm}
  \caption{One implicit Euler step: the update $x_{k+1}$ lies on the segment
    between $x_k$ and the minimizer $x^\star$}
  \label{fig:flow}
\end{figure}

\section{Analysis}
\label{sec:analysis}

The workhorse is the one-step descent estimate for the iterates of
\cref{eq:implicit} defined in \cref{sec:method}.

\begin{lemma}
  \label{lem:descent}
  For every $k \ge 0$,
  \begin{equation}
    \label{eq:descent}
    f(x_{k+1}) - f^\star \le \frac{1}{2\tau}
      \bigl(\lVert x_k - x^\star \rVert^2
      - \lVert x_{k+1} - x^\star \rVert^2\bigr).
  \end{equation}
\end{lemma}

Strong convexity converts \cref{eq:descent} into a contraction, giving the
main result: with $\tau = 1/L$, \cref{thm:linear} below bounds the error of
the scheme from \cref{sec:method} whose geometry \cref{fig:flow} depicts.

\begin{theorem}
  \label{thm:linear}
  The iterates of \cref{eq:implicit} satisfy
  $\lVert x_k - x^\star \rVert^2 \le (1 + \mu/L)^{-k}
  \lVert x_0 - x^\star \rVert^2$.
\end{theorem}

\Cref{lem:descent} and \cref{thm:linear} together close the argument opened
in \cref{sec:intro}: the discrete flow inherits the exponential stability of
its continuous parent, with constants that depend only on $\mu$ and $L$.

\end{document}
